\section{Dokumentacja testów}

\subsection{Cel i zakres testowania}

Celem testów jest sprawdzenie reguł systemu (personalizacji, adaptacji, oceniania) oraz uprawnień i poprawności wszystkich endpointów serwera. Testy obejmują wszystkie moduły serwera: konta, kursy i klasy, treść lekcji i pliki, model ucznia oraz quizy, a także dane demonstracyjne.

\subsection{Strategia testowania}

\begin{center}
\begin{tabular}{lll}
\toprule
\textbf{Poziom} & Jednostkowy & Integracyjny \\
\midrule
\textbf{Typ} & white-box & black-box HTTP \\
\textbf{Narzędzie} & pytest & pytest + FastAPI TestClient \\
\textbf{Zakres} & czyste funkcje reguł & endpointy REST z bazą danych \\
\bottomrule
\end{tabular}
\end{center}

\bigskip

\textbf{Testy jednostkowe} sprawdzają czyste funkcje (punktacja ankiety, zwijanie materiałów, styl efektywny, aktualizacja wag, propozycja stylu, ocenianie quizu) -- bez bazy danych i bez frameworka.

\textbf{Testy integracyjne} wysyłają żądania HTTP do aplikacji przez \texttt{TestClient} i korzystają z prawdziwej bazy PostgreSQL (osobna baza \texttt{imsu\_test}). Schemat bazy testowej jest budowany tymi samymi migracjami Alembic co baza aplikacji, a każdy test działa w transakcji wycofywanej po jego zakończeniu, więc testy są od siebie niezależne. Wspólne dane (uczeń, nauczyciel, administrator, kurs z lekcjami) dostarczają fixture'y pytest.

Interfejs użytkownika sprawdzany jest ręcznie według scenariusza na kontach demonstracyjnych; kod klienta przechodzi lint (ESLint), kontrolę formatowania (Prettier) i build produkcyjny.

\subsection{Środowisko testowe}

\begin{center}
\begin{tabular}{ll}
\toprule
\textbf{Narzędzie} & \textbf{Wersja / uwagi} \\
\midrule
Python & 3.12+ \\
pytest & testy jednostkowe i integracyjne \\
FastAPI TestClient (httpx) & żądania HTTP do aplikacji \\
PostgreSQL & 17, baza \texttt{imsu\_test} w kontenerze Docker \\
Alembic & budowa schematu bazy testowej \\
Ruff & lint i formatowanie kodu serwera \\
\bottomrule
\end{tabular}
\end{center}

\subsection{Przypadki testowe}

Poniższa tabela zawiera reprezentatywne przypadki testowe. Kolumna \emph{Typ}: \textbf{J} -- jednostkowy, \textbf{I} -- integracyjny (HTTP).

\begingroup
\scriptsize
\setlength{\tabcolsep}{2pt}
\begin{longtable}{p{1cm} p{3.2cm} p{1.6cm} p{0.6cm} p{4.3cm} p{4.3cm}}
\toprule
\textbf{ID} & \textbf{Nazwa} & \textbf{Moduł} & \textbf{Typ} & \textbf{Dane wejściowe / warunki} & \textbf{Oczekiwany wynik} \\
\midrule
\endfirsthead
\toprule
\textbf{ID} & \textbf{Nazwa} & \textbf{Moduł} & \textbf{Typ} & \textbf{Dane wejściowe / warunki} & \textbf{Oczekiwany wynik} \\
\midrule
\endhead
\midrule
\multicolumn{6}{r}{\footnotesize Kontynuacja na następnej stronie\ldots} \\
\endfoot
\bottomrule
\endlastfoot

TC-01 & Rejestracja ucznia & identity & I &
  \texttt{POST /auth/register}: rola \emph{student} &
  201; ciasteczko sesji \texttt{HttpOnly} \faCheck \\
\addlinespace
TC-02 & Rejestracja jako administrator & identity & I &
  \texttt{POST /auth/register}: rola \emph{admin} &
  422; konto nie powstaje \faCheck \\
\addlinespace
TC-03 & Blokada konta działa od razu & identity & I &
  administrator blokuje zalogowanego użytkownika &
  kolejne żądanie \texttt{GET /auth/me}: 401 \faCheck \\
\addlinespace
\midrule
TC-04 & Katalog kursów ucznia & curriculum & I &
  \texttt{GET /courses}: kursy prywatne i opublikowane &
  uczeń widzi tylko kursy opublikowane \faCheck \\
\addlinespace
TC-05 & Kurs prywatny przez klasę & curriculum & I &
  kurs prywatny przypisany do klasy ucznia &
  uczeń otwiera kurs, lekcję i quiz; po usunięciu z klasy -- 404 \faCheck \\
\addlinespace
TC-06 & Cudzy kurs & curriculum & I &
  nauczyciel otwiera kurs innego nauczyciela &
  404 (bez zdradzania istnienia kursu) \faCheck \\
\addlinespace
\midrule
TC-07 & Plik niedozwolonego typu & content & I &
  \texttt{POST /files}: plik \texttt{.svg} lub \texttt{.exe} &
  400 \texttt{FILE\_TYPE\_NOT\_ALLOWED} \faCheck \\
\addlinespace
TC-08 & Plik innego nauczyciela w lekcji & content & I &
  \texttt{PUT /lessons/\{id\}/blocks}: załącznik z cudzym plikiem &
  400 \texttt{FILE\_NOT\_FOUND} \faCheck \\
\addlinespace
\midrule
TC-09 & Ankieta bez klucza & learning\_style & I &
  \texttt{GET /survey/questions} &
  odpowiedzi bez informacji o formie V/A/K \faCheck \\
\addlinespace
TC-10 & Remis w ankiecie & learning\_style & I &
  tyle samo punktów V i A &
  styl mieszany \texttt{VA} \faCheck \\
\addlinespace
TC-11 & Zwinięte materiały & learning\_style & I &
  uczeń o stylu V otwiera lekcję z materiałami A i K &
  wszystkie bloki; materiały A i K na liście zwiniętych \faCheck \\
\addlinespace
TC-12 & Uczeń bez ankiety & learning\_style & I &
  \texttt{GET /lessons/\{id\}/personalized-blocks} &
  403 \texttt{SURVEY\_REQUIRED} \faCheck \\
\addlinespace
TC-13 & Pierwszy sukces i pierwsza porażka & learning\_style & J &
  kolejne wyniki quizu w tej samej lekcji &
  liczy się tylko pierwszy sukces i pierwsza porażka; powtórzenia nie zmieniają wag \faCheck \\
\addlinespace
TC-14 & Propozycja zmiany stylu & learning\_style & J &
  wagi V 6 / A 9 / K 2, styl V &
  propozycja: A \faCheck \\
\addlinespace
TC-15 & Styl ręczny wyłącza propozycje & learning\_style & J &
  jak TC-14, źródło stylu: ręczne &
  brak propozycji \faCheck \\
\addlinespace
\midrule
TC-16 & Pytania według stylu & assessment & I &
  uczeń o stylu V pobiera quiz &
  pytania wspólne i V, bez klucza odpowiedzi \faCheck \\
\addlinespace
TC-17 & Rozwinięty materiał w quizie & assessment & I &
  uczeń rozwija materiał innej formy i pobiera quiz ponownie &
  w quizie pojawia się pytanie o ten materiał \faCheck \\
\addlinespace
TC-18 & Wielokrotne podejścia & assessment & I &
  dwa zaliczone podejścia w tej samej lekcji &
  wagi zmieniają się tylko po pierwszym \faCheck \\
\addlinespace
TC-19 & Pytanie spoza stylu & assessment & I &
  odpowiedź na pytanie formy, której uczeń nie widzi &
  400 \texttt{QUIZ\_ANSWER\_MISMATCH} \faCheck \\
\addlinespace
\midrule
TC-20 & Idempotentny seed & demo & I &
  dwukrotne uruchomienie seeda &
  brak duplikatów kursów, lekcji i klasy \faCheck \\

\end{longtable}
\endgroup

\subsection{Scenariusze testowe}

\subsubsection{ST-01: Od ankiety do adaptacji}
\textbf{Cel:} sprawdzenie pełnej ścieżki ucznia.
\begin{enumerate}[leftmargin=*]
    \item Uczeń bez ankiety próbuje otworzyć lekcję (TC-12) -- 403 \texttt{SURVEY\_REQUIRED}.
    \item Uczeń wypełnia ankietę -- powstaje profil ze stylem i wagami (TC-09, TC-10).
    \item Uczeń otwiera lekcję -- materiały innych stylów są zwinięte (TC-11).
    \item Uczeń rozwija materiał i rozwiązuje quiz -- pytania odpowiadają stylowi efektywnemu (TC-16, TC-17).
    \item Wynik quizu zmienia wagi; kolejne podejście już ich nie zmienia (TC-18).
\end{enumerate}

\subsubsection{ST-02: Propozycja zmiany stylu}
\textbf{Cel:} sprawdzenie reakcji systemu na wyniki ucznia.
\begin{enumerate}[leftmargin=*]
    \item Wagi ucznia o stylu V spadają po niezaliczonych quizach do V 7 / A 8 / K 3.
    \item Profil ucznia zawiera propozycję zmiany stylu na A z uzasadnieniem (TC-14).
    \item Uczeń odrzuca propozycję -- nie pojawia się ona ponownie.
    \item Uczeń wybiera styl ręcznie -- propozycje są wyłączone (TC-15).
\end{enumerate}

\subsubsection{ST-03: Dostęp do kursów przez klasy}
\textbf{Cel:} sprawdzenie widoczności kursów prywatnych.
\begin{enumerate}[leftmargin=*]
    \item Nauczyciel tworzy klasę, dodaje ucznia i przypisuje kurs prywatny.
    \item Uczeń nie widzi kursu w katalogu (TC-04), ale otwiera go przez klasę (TC-05).
    \item Uczeń spoza klasy i inny nauczyciel otrzymują 404 (TC-06).
    \item Po usunięciu ucznia z klasy kurs przestaje być dla niego dostępny.
\end{enumerate}

\subsection{Podsumowanie pokrycia testami}

\begin{center}
\begin{tabular}{lrrr}
\toprule
\textbf{Moduł} & \textbf{Jednostkowe} & \textbf{Integracyjne} & \textbf{Łącznie} \\
\midrule
\texttt{identity}        & 6  & 42  & 48 \\
\texttt{curriculum}      & -- & 64  & 64 \\
\texttt{content}         & 2  & 23  & 25 \\
\texttt{learning\_style} & 23 & 42  & 65 \\
\texttt{assessment}      & 6  & 15  & 21 \\
\texttt{demo}            & -- & 10  & 10 \\
\texttt{core} (błędy, stan serwera) & -- & 5 & 5 \\
\midrule
\textbf{Razem}           & \textbf{37} & \textbf{201} & \textbf{238} \\
\bottomrule
\end{tabular}
\end{center}

Tabela podaje liczbę funkcji testowych. Część testów jest sparametryzowana (ten sam test dla wielu danych wejściowych), dlatego pytest wykonuje łącznie \textbf{349} przypadków testowych.